\documentclass[11pt]{article}

\usepackage[utf8]{inputenc}
\usepackage[T1]{fontenc}
\usepackage{geometry}
\usepackage{amsmath}
\usepackage{amssymb}
\usepackage{booktabs}
\usepackage{hyperref}
\usepackage{xurl}
\usepackage{xcolor}
\usepackage{listings}
\usepackage{enumitem}
\usepackage{parskip}
\usepackage{graphicx}
\graphicspath{{../images/}}

\geometry{margin=1in}

\lstset{
  language=Java,
  basicstyle=\ttfamily\small,
  keywordstyle=\color{blue},
  commentstyle=\color{gray},
  stringstyle=\color{teal},
  showstringspaces=false,
  columns=fullflexible,
  frame=single,
  framerule=0.2pt,
  breaklines=true
}

\setlist[itemize]{topsep=0.3em,itemsep=0.2em}
\setlist[enumerate]{topsep=0.3em,itemsep=0.2em}

% Simple per-section source line.
\newcommand{\notesources}[1]{\par\noindent{\small\color{gray}\textit{Sources:} \sloppy #1\par}}

% Unit-wise source strings for this subject (used by slides/notes).
% Path is relative to the lecture `latex/` folder (the compilation CWD).
% OOPS with Java (CSEG1044) — unit-wise sources
%
% These are short, student-facing reference strings intended to be used with:
%   - \slidesources{...}  (in beamer slides)
%   - \notesources{...}   (in lecture notes)
%
% Style inspiration: other subjects in My-Subjects/ use semicolon-separated sources
% and include an "accessed" date for web documentation.

\newcommand{\OOPSUnitISources}{Schildt, \textit{Java: A Beginner's Guide} (10e), McGraw-Hill (Java fundamentals + OOP basics).; Oracle Java Tutorials: Getting Started + Language Basics + Classes and Objects (accessed \today).; Java SE API docs (e.g., \texttt{String}, \texttt{Scanner}) (accessed \today).}

\newcommand{\OOPSUnitIISources}{Schildt, \textit{Java: The Complete Reference} (12e), McGraw-Hill (classes, inheritance, packages, interfaces).; Bloch, \textit{Effective Java} (3e), Addison-Wesley, 2018 (design choices; interfaces vs abstract classes).; Oracle Java Tutorials: Classes and Objects + Interfaces + Packages (accessed \today).; Java SE API docs (e.g., \texttt{Comparable}, \texttt{Runnable}) (accessed \today).}

\newcommand{\OOPSUnitIIISources}{Schildt, \textit{Java: The Complete Reference} (12e) (nested/inner classes, exceptions, threads, I/O).; Oracle Java Tutorials: Nested Classes + Exceptions + Concurrency + Basic I/O (accessed \today).; Java SE API docs (e.g., \texttt{Thread}, \texttt{AutoCloseable}) (accessed \today).}

\newcommand{\OOPSUnitIVSources}{Schildt, \textit{Java: The Complete Reference} (12e) (generics, lambdas, Swing, JDBC).; Bloch, \textit{Effective Java} (3e) (generics + lambdas).; Oracle Java Tutorials: Generics + Lambda Expressions + Swing + JDBC Basics (accessed \today).; Java SE API docs (e.g., \texttt{java.util.function}, \texttt{javax.swing}, \texttt{java.sql}) (accessed \today).}

\newcommand{\OOPSUnitVSources}{Schildt, \textit{Java: The Complete Reference} (12e) (Collections Framework + wrapper classes).; Oracle Java docs: Collections Framework overview (accessed \today).; Java SE API docs (\texttt{java.util} collections; wrapper classes in \texttt{java.lang}) (accessed \today).}

\newcommand{\OOPSUnitVISources}{Git SCM: \textit{Pro Git} (2e) (version control workflow), accessed \today.; GitHub Docs (repositories, issues, pull requests), accessed \today.; Oracle Java Tutorials: Swing + JDBC (accessed \today).; JUnit 5 User Guide (accessed \today).; Apache Maven Project: Getting Started (accessed \today).; Gradle User Manual: Getting Started (accessed \today).; UML class diagrams: UML 2.x overview/spec (accessed \today).}

\newcommand{\OOPSMiscWhyInterfacesSources}{Schildt, \textit{Java: The Complete Reference} (12e) (interfaces).; Bloch, \textit{Effective Java} (3e) (prefer interfaces where appropriate).; Oracle Java Tutorials: Interfaces (accessed \today).; Java SE API docs (e.g., \texttt{Comparable}, \texttt{Runnable}) (accessed \today).}



\title{Object Oriented Programming (CSEG1044)\\Unit 02 -- Lecture 04 Notes\\Inheritance, `super`, Overriding, and Dynamic Dispatch}
\begin{document}
\maketitle
\tableofcontents

\section{Lecture Snapshot}
\begin{itemize}
  \item \textbf{Unit focus:} Classes, Inheritance, Packages, and Interfaces
  \item \textbf{Main new ideas:} Inheritance and types (conceptual), \texttt{super} usage, Overriding and dynamic dispatch
  \item \textbf{Course thread:} Use Inheritance, `super`, Overriding, and Dynamic Dispatch to strengthen the domain model, APIs, and access rules inside Campus Resource Manager.
\end{itemize}
\notesources{\OOPSUnitIISources}

\section{Lecture Overview}
Inheritance and polymorphism are two core tools in OOP, but they must be used carefully.
This lecture focuses on:
\begin{itemize}
  \item modeling \textbf{is-a} relationships with inheritance,
  \item using \texttt{super} correctly,
  \item overriding methods and understanding \textbf{dynamic dispatch}.
\end{itemize}
\notesources{\OOPSUnitIISources}

\section{Prerequisite Bridge}
Before reading the new material in depth, do a fast recall of the older ideas listed below.
\begin{itemize}
  \item \textbf{Class design}: Inheritance helps only when the parent and child already have meaningful responsibilities. Main reference: \texttt{Unit 02 / Lecture 01 slides/notes}.
  \item \textbf{Constructors}: The rule for super(...) depends on the constructor logic from the previous lecture. Main reference: \texttt{Unit 02 / Lecture 02 slides/notes}.
\end{itemize}
This lecture only revisits those ideas briefly so that the main explanation can stay focused on the new concept sequence.
\notesources{\OOPSUnitIISources}

\section{Learning Outcomes}
After this lecture, you should be able to:
\begin{itemize}
  \item Decide is-a (inheritance) vs has-a (composition).
  \item Use \texttt{super(...)} and \texttt{super.method()} correctly.
  \item Override methods and explain dynamic dispatch with an example.
\end{itemize}

\section{Suggested Reading Order}
\begin{enumerate}
  \item Read the \textbf{Prerequisite Bridge} first and confirm each recalled idea.
  \item Study the central explanation in this order: \textit{Inheritance and types (conceptual)} $\rightarrow$ \textit{\texttt{super} usage} $\rightarrow$ \textit{Overriding and dynamic dispatch}.
  \item Trace the worked example line by line before checking short answers.
  \item Attempt the practice section without looking at the solution immediately.
  \item Finish with the exit questions and further-study suggestions to confirm retention.
\end{enumerate}
\notesources{\OOPSUnitIISources}

\section{Key Terms (Quick)}
\begin{itemize}
  \item \textbf{Inheritance}: a subclass extends a superclass (\texttt{extends}).
  \item \textbf{Superclass/Subclass}: parent/child in inheritance hierarchy.
  \item \textbf{Overriding}: subclass provides its own implementation of a superclass method (same signature).
  \item \textbf{Dynamic dispatch}: method call resolved at runtime based on the actual object type.
  \item \textbf{Upcasting}: treating a subclass object as a superclass reference, for example \path|Shape s = new Circle(...)|.
  \item \textbf{Composition}: has-a relationship (\texttt{Car} has an \texttt{Engine}).
\end{itemize}

\section{Detailed Notes}
\subsection{Inheritance and types (conceptual)}
Java supports:
\begin{itemize}
  \item single inheritance of implementation (one direct superclass),
  \item multilevel and hierarchical inheritance patterns,
  \item multiple inheritance of \emph{type} through interfaces (covered more later).
\end{itemize}

\textbf{Is-a vs has-a:}
\begin{itemize}
  \item Use inheritance only when the subclass is truly a kind of superclass.
  \item If something is a part of another thing, prefer composition.
\end{itemize}
Examples:
\begin{itemize}
  \item Dog is-a Animal (inheritance)
  \item Car has-a Engine (composition)
\end{itemize}

\subsection{\texttt{super} usage}
\textbf{Constructor:} a subclass constructor must call a superclass constructor.
If you do not write \texttt{super(...)}, Java inserts \texttt{super()} automatically \emph{only if} the superclass has a no-arg constructor.

\textbf{Rule:} \texttt{super(...)} must be the \textbf{first statement} in the constructor (similar to \texttt{this(...)}).

\textbf{Member access:} use \texttt{super.method()} to explicitly call the parent method if overridden.

\subsection{Overriding and dynamic dispatch}
Overriding rules (important ones):
\begin{itemize}
  \item same method signature in subclass,
  \item cannot reduce visibility (\texttt{public} cannot become \texttt{protected/private}),
  \item use \texttt{@Override} to catch mistakes.
\end{itemize}

\textbf{Dynamic dispatch:} if you store a circle as \path|Shape s = new Circle(...)|, then calling \texttt{s.area()} still runs \texttt{Circle.area()} at runtime.
This is the heart of runtime polymorphism.
\notesources{\OOPSUnitIISources}

\section{Running Project Connection}
Use Inheritance, `super`, Overriding, and Dynamic Dispatch to strengthen the domain model, APIs, and access rules inside Campus Resource Manager.
\begin{itemize}
  \item Use Campus Resource Manager as the running case study, or map the same idea to your own project domain if you are using another example.
  \item Ask where today's idea should live: model, service, DAO, UI, or team workflow.
  \item Use the lecture example as a smaller version of the same design choice you will make in the capstone.
\end{itemize}
\notesources{\OOPSUnitIISources}

\section{Common Mistakes and Confusions}
\begin{itemize}
  \item Choosing inheritance just to reduce code repetition quickly.
  \item Forgetting that super(...) must run before subclass-specific initialization.
  \item Confusing overriding with overloading.
\end{itemize}
\notesources{\OOPSUnitIISources}

\section{Worked Example: \texttt{Shape} hierarchy}
\begin{lstlisting}
class Shape {
    private final String color;

    Shape(String color) {
        this.color = color;
    }

    double area() { return 0.0; }
    String getColor() { return color; }
}

class Circle extends Shape {
    private final double r;

    Circle(String color, double r) {
        super(color); // superclass constructor
        this.r = r;
    }

    @Override
    double area() { return Math.PI * r * r; }
}
\end{lstlisting}

\begin{lstlisting}
public class Demo {
    public static void main(String[] args) {
        Shape s = new Circle("red", 2.0); // upcasting
        System.out.println(s.area());    // dynamic dispatch
    }
}
\end{lstlisting}

\section{Demo / Observation Task}
Use this block as a short reproducible demo or observation task.
\begin{itemize}
  \item Reproduce the lecture's main example around \textit{Inheritance and types (conceptual)} once without looking at the solution first.
  \item Change one input, rule, or design choice in Campus Resource Manager and predict the result before checking it.
  \item Record one success case, one edge case, and one failure case so the idea becomes observable instead of purely theoretical.
\end{itemize}
\notesources{\OOPSUnitIISources}

\section{Practice Check (with brief solutions)}
\begin{enumerate}
  \item Is-a or has-a?\; \textbf{(a)} Student--IDCard \textbf{(b)} Manager--Employee\par
  \textbf{Answer:} (a) has-a, (b) is-a (in a simple model).

  \item What does \texttt{super(...)} do in a constructor?\par
  \textbf{Answer:} calls the superclass constructor to initialize the parent part of the object.

  \item If \texttt{Shape s = new Circle(...)} and you call \texttt{s.area()}, which implementation runs?\par
  \textbf{Answer:} \texttt{Circle.area()} (runtime dispatch).
\end{enumerate}

\section{Self-Study Practice (no solutions)}
\begin{enumerate}
  \item Create \texttt{Rectangle} and \texttt{Triangle} subclasses of \texttt{Shape} and override \texttt{area()}.
  \item Design a model using composition (has-a) and justify why inheritance is not appropriate.
  \item Write one example where you call \texttt{super.method()} from an overridden method.
\end{enumerate}

\section{Further Study}
\begin{itemize}
  \item Revisit the prerequisite references if any recall bullet still feels weak.
  \item Rework the lecture example with one changed assumption, input, or design choice.
  \item Read the textbook and documentation references listed in the source line for a second explanation of the same topic.
  \item Apply the same idea once inside Campus Resource Manager so that the concept moves from theory to design practice.
\end{itemize}
\notesources{\OOPSUnitIISources}

\section{Exit Questions (with sample answers)}
\textbf{Q1.} Why is inheritance sometimes risky?\par
\textbf{Sample answer:} Inheritance creates strong coupling between parent and child. If the parent changes, many subclasses may break. Using inheritance only when the relationship is truly is-a helps avoid bad designs.

\vspace{0.6em}
\textbf{Q2.} What is dynamic dispatch?\par
\textbf{Sample answer:} Dynamic dispatch is when the method that executes is chosen at runtime based on the actual object type, not the reference type.

\vspace{0.6em}
\textbf{Q3.} How does dynamic dispatch support extensible design?\par
\textbf{Sample answer:} Client code can depend on a parent type while the correct overridden behavior is selected at runtime.

\end{document}
